\chapter{Implementierung}
\label{chap:implementierung}

\section{Ausgangszustand der Check-API und Aufgabenverarbeitung}
Damit verständlich wird, wie die Änderungen durch diese Bachelorarbeit das System verändern, gehe ich in diesem Abschnitt darauf ein, wie das System im nahen Umfeld der Änderungen funktioniert und Aufgebaut ist.

\subsection{Übungsaufgabenverwaltung}
In diesem Unterabschnitt gehe ich genauer auf die Bereitstellung, Verarbeitung und Speicherung von Übungsaufgaben ein.

\subsubsection{Koan-Repository}
Das Koan-Repository ist ein Git-Repository, in welchem sämtliche Übungsaufgaben abgelegt sind. Diese Aufgaben können von Nutzern mit der Rolle Professor über das Frontend geladen werden. Dabei wird mit spezieller Berechtigung das Repository angefragt und CodeTask erhält die Aufgaben in der Haupt-API.

\subsubsection{Koan-Parser}
Die Aufgaben werden daraufhin in speziell festgelegte Elemente zerteilt. Dafür zuständig ist der Parser, dieser besteht aus Reg-ex Regeln, welche ausdrücken, wo ein Element Anfängt, Aufhört und was dazwischen relevant ist.

\subsubsection{Datenbank}
Die aufgeteilten Aufgaben werden daraufhin in der Datenbank gespeichert. Das sorgt dafür, das die Aufgaben nicht für jede Anfrage aus dem Repository geladen und durch den Parser verarbeitet werden müssen.

\subsubsection{Frontend}
Öffnet ein Nutzer eine Lektion, und wählt darin eine Aufgabe, welche er lösen möchte, so wird von der Haupt-API die korrekte Aufgabe geladen und mit allen Informationen im Frontend dar gestellt.

\section{Erweiterung des Parsers und der Testdatenverarbeitung}
Angefangen mit der Implementierung der Infrastruktur für die versteckten Tests, mussten Änderungen an dem Parser vorgenommen werden. Für eine saubere Aufteilung musste erstmals verstanden werden, wie der Parser bisher gearbeitet hat. In dem Listing~\ref{lst:koanExampleWithoutHiddenTests} wird die CodeTask Aufgabe von einem Koan abgebildet. Dieser Aufgabenblock ist speziell Formatiert und nutzt spezielle Markierungen wie \texttt{//solve}, \texttt{//endsolve}, \texttt{//test} und \texttt{//endtest}. Mithilfe dieser Markierungen kann der Parser gezielt nach diesen Bereichen suchen und diese an diesen Stellen trennen.
Alles was vor dem \texttt{//solve} steht, dient als Hilfestellung für den Nutzer und wird ohne weiteres übernommen. Der Inhalt in dem  \texttt{solve} Block ist die Musterlösung für diese Aufgabe und muss getrennt werden, damit diese nicht von Anfang an dem Nutzer angezeigt wird. In dem \texttt{test} Block ist der öffentliche Test definiert und soll dem Nutzer ebenfalls als Hilfestellung angezeigt werden.
Der Parser hat aus diesen Blöcken \texttt{code}, und \texttt{codeSolution} gemacht. Hier fehlen aber die Tests. Das Feld \texttt{code} wurde im Parser wie im Listing~\ref{lst:codeFieldOldParser} gezeigt zusammen gesetzt. Dabei wird der erste Bereich übernommen, der \texttt{solve} Block entfernt und durch ein \texttt{\_\#\_} ersetzt und zum Schluss wird der \texttt{test} Block wieder angehängt. Das Feld \texttt{codeSolution} besteht aus dem \texttt{solve} Block.
Diese Lösung ist im Prinzip nicht schlecht, reicht für versteckte Tests aber nicht mehr aus.
Anstelle das \texttt{code} Feld auch die Tests enthält, wird jetzt nur noch der Code vor dem \texttt{//solve} zusammen mit dem \texttt{\_\#\_} gespeichert. Neu dazu kommen die Felder \texttt{tests} und \texttt{hiddenTests}. Der Parser speichert also jeden Block einzeln ab. Das erleichtert bzw. ermöglicht die spätere Bereitstellung der relevanten Daten.
Der Parser wurde so erweitert, das er die Möglichkeit hat, versteckte Tests zu erkennen. Die Erweiterung, des in Listing~\ref{lst:koanExampleWithoutHiddenTests} gezeigten Koans, kann in Listing~\ref{lst:koanExampleWithHiddenTests} gesehen werden. Koans können um einen \texttt{//hiddentest} und \texttt{//enhiddentest} Block erweitert werden und können ohne weiter Anpassung des Parsers verarbeitet werden.

\begin{lstlisting}[
	language=scala,
	caption={Koan Ausschnitt ohne versteckte Tests}, label={lst:koanExampleWithoutHiddenTests},
	float=htbp
	]
codetask(
"""# Übung: Mehrzeilige Adresse mit Interpolation

Nutze einen mehrzeiligen interpolierten String mit `stripMargin`, um eine Adresse aufzubauen.
"""
) {
  // val name = "HIERÄNDERN"
  // val street = "123 Navy St"
  // val city = "HIERÄNDERN"
  // val address =
  // s"""
  //     |$HIERÄNDERN
  //     |$HIERÄNDERN
  //     |$HIERÄNDERN
  //""".stripMargin
  //solve
  val name = "Grace Hopper"
  val street = "123 Navy St"
  val city = "Arlington, VA 22201"
  val address = s"""
    |$name
    |$street
    |$city
  """.stripMargin
  //endsolve
	
  //test
  address should include ("Grace Hopper")
  address should include ("123 Navy St")
  address should include ("Arlington")
  //endtest
}
\end{lstlisting}
\begin{lstlisting}[
	language=scala,
	caption={Das Feld "code" nach dem Bestehenden Parser}, label={lst:codeFieldOldParser},
	float=htbp
	]
code= 
"""// val name = "HIERÄNDERN"
// val street = "123 Navy St"
// val city = "HIERÄNDERN"
// val address =
// s"""
//     |$HIERÄNDERN
//     |$HIERÄNDERN
//     |$HIERÄNDERN
//""".stripMargin
_#_
address should include ("Grace Hopper")
address should include ("123 Navy St")
address should include ("Arlington")"""
\end{lstlisting}
\begin{lstlisting}[
	language=scala,
	caption={Koan Ausschnitt mit versteckten Tests}, label={lst:koanExampleWithHiddenTests},
	float=htbp
	]
...
//test
address should include ("Grace Hopper")
address should include ("123 Navy St")
address should include ("Arlington")
//endtest

//hiddentest
address.contains("\n") shouldBe true
val lines = address.linesIterator.toList
lines.length shouldBe 3
lines.mkString(" ") should include (name)
lines.mkString(" ") should include (street)
lines.mkString(" ") should include (city)
//endhiddentest
\end{lstlisting}
	
	
\section{Bereitstellung und Zuordnung von versteckten Tests}
Wie gerade schon angesprochen werden mit den Veränderungen im Parser die Blöcke in einzelne Felder abgelegt und über die bestehende Datenbankverbindung in die Datenbank geschrieben. Für die Speicherung in der Datenbank mussten lediglich die neuen Felder übergeben werden. Die beschriebene Änderung erlaubt dem System eine getrennte Übertragung der einzelnen Blöcke. Die Musterlösung und der Code, welcher die Aufgabenstellung helfend unterstützt muss zusammen mit den öffentlichen Tests über die Haupt-API an das Frontend geschickt werden. Hier muss beachtet werden, das die versteckten Tests nicht ebenfalls mitgeschickt werden, da diese sonst ausgelesen werden können. Löst der Nutzer die Aufgabe und schickt sein Ergebnis ab, so wird zusammen mit dem Nutzercode, die öffentlichen Tests, die Lösung und Informationen über die Aufgabe selbst mitgeschickt. Diese Informationen werden in der Haupt-API genutzt um die korrekten versteckten Tests in der Datenbank zu finden, sodass der Nutzercode zusammen mit den öffentlichen und versteckten Tests an die Check-API übergeben werden kann. 

\section{Anpassung der Check-API}
Unter die Anpassungen der Check-API fallen auf die Erweiterung der Schnittstelle zur Unterstützung von versteckten Tests, eine erhöhte Sicherheit und Isolation, Verbesserungen der Leistung sowie eine verbesserte Rückmeldung für den Nutzer.
Die Check-API muss die Anfrage von der Haupt-API mit den neuen Änderungen auch entsprechend entgegen nehmen. Da von der Haupt-API die Felder \texttt{code}, \texttt{tests} und \texttt{hiddenTests} übergeben werden, müssen diese auch in der Check-API angenommen und verarbeitet werden. In der Check-API angekommen, werden die Felder an den sogenannten codeTester übergeben. Hier werden die drei Felder erstmals wieder zu einem großen vereinten Stück Code in Form von einem String. Dieser String wird mit einer Liste voller verbotener Schlüsselwörter geprüft. Sollte auch nur eines dieser Schlüsselwörter in dem Code vorkommen, wird die Ausfürhung sofort mit einem Failure und einer Exception abgebrochen. Die Prüfung verhindert ausgewählte bekannte Problemfälle, ersetzt jedoch keine vollständige Sicherheitsabsicherung der Codeausführung. Darunter z.~B. \texttt{System.exit}. Dieser Befehl, schaltet den gesamten Prozess aus in dem er läuft. Nach dieser Sicherheitsüberprüfung werden die einzelnen Blöcke leicht formatiert und in ein vordefiniertes Template eingesetzt. Dieses Template beinhaltet alle notwendigen Imports, sowie Methoden um die Tests korrekt ausführen zu können. Das befüllte Template ist ab diesem Zeitpunkt vollständig un bereit zur Compilierung. Dafür wird das Template in eine virtuelle Datei geschrieben. Eine virtuelle Datei existiert dabei nur im Arbeitsspeicher und wird nicht auf die Festplatte geschrieben. Das ermögliche schnellere Geschwindigkeiten. Die Idee ist nicht neu sonder war bereits in dem Vorhandenen Compiler eingebaut. Ist diese virtuelle Datei erstellt kann der Compiler gestartet werden. Der Compiler übersetzt den Code, sodass der Code ausgeführt werden kann. Der übersetzte Code wird temporär auf die Festplatte geschrieben. Das bringt den Vorteil, das neuerdings die Codeausführung mehrerer Anfragen parallel ausgeführt werden können und so eine schnellere Antwortzeit auch unter Last ermöglicht wird.. Die Parallelität ist dabei allerdings von der Hardware begrenzt. Würden pro Anfrage dauerhaft, ungebremst neue Unterprozesse erstellt werden, würde das System sich selbst verlangsamen, da mit der Erstellung eines Unterprozesses ein gewisser aufwand verbunden ist. Daher begrenzen wir die Anzahl der gleichzeitig existierenden Unterprozesse. Mindestens wird ein Unterprozess erstellt, maximal aber $\texttt{Anzahl} = \frac{\text{CPU-Kerne}}{2}$. 
Das führt dazu, das wir nicht die gesamte Leistung beanspruchen und noch genügend Leistung für das restliche System besteht.
In dem neu erstellten Unterprozess wird der auf der Festplatte gespeicherte compilierte Code ausgeführt. Diese Dateien müssen nach dem Ausführen zum Schluss noch aufgeräumt werden. Damit ist der grobe Ablauf und die Implementierung in Teilen beschrieben. Etwas tiefer wird auf die Leistungsoptimierung noch in Kapitel \ref{chap:evaluation} Evaluation eingegangen.

\section{Zeitüberschreitung}
Die Check-API wurde noch in weiteren Bereichen verbessert, darunter das Handhaben von Zeitüberschreitungen. Direkt nach dem der Unterprozess, zuständig für das ausführen des Codes, gestartet ist, wird ein Thread gestartet, dessen einziger Job es ist, den Unterprozess zu überwachen. Der Thread wartet blockierend, bis der Unterprozess beendet ist. Das passiert entweder durch einen Laufzeitfehler oder durch das erfolgreiche erreichen des Ende vom Code. Solange der Thread blockierend wartet, wartet ebenfalls der Hauptprozess, entweder darauf, das der Thread ein Ergebnis bekommen hat oder eine gewisse Zeit abgelaufen ist. Aktuell 5 Sekunden, dies ist allerdings frei anpassbar, sollte dieses Zeitfenster zu eng gewählt worden sein. Tritt der Fall ein, das der Hauptprozess diese Zeit vollständig warten musste, greift er ein und beendet den Unterprozess der gerade den Code ausführt. Dieses Verhalten sorgt dafür, das fehlerhafter Code nicht unendlich lange laufen kann und permanent Ressourcen verbraucht. Solch eine lange Laufzeit kann mehrere Ursachen haben. Darunter eine klassische unendliche Schleife, welche sich für immer wiederholt. Das Ergebnis ist, der Code wird für eine unendliche Dauer ausgeführt. Eine zweite Ursache kann unoptimierter Code sein, der sehr viele unterschiedliche Fälle durchlaufen muss und so sehr lange in der Ausführung braucht. Dieser Fall kann problematisch sein, da dieser Code eventuell sogar komplett korrekt sein kann und dennoch abgebrochen wird. Im Fall von CodeTask stellt das allerdings kein besonderes Problem da, die Aufgaben zum lernen recht klein gewählt werden und keine extrem komplexen Algorithmen mit tausenden Daten fordern. Sollte es dennoch zu diesem Problem kommen, kann das Zeitlimit angepasst werden.

\section{Fehlerbehandlung und Strukturierte Ergebnisrückgabe}
Im Vergleich zum ursprünglichen System, wurde die Check-API ebenfalls um einige Fehlerbehandlungen und Ergebnisrückgaben erweitert. Ursprünglich hat die Check-API lediglich mit zwei unterschiedlichen Fällen geantwortet. Einmal mit einem \texttt{400 Bad Request}  und einmal mit einem \texttt{200 Ok} Statuscode. Bei einem \texttt{400 Bad Request} wurde dabei eine Nachricht mitgeschickt, welche vermuten lässt, das versucht wurde ein Compiler oder Laufzeitfehler mitschicken zu können, hier ist allerdings etwas schief gegangen, sodass lediglich immer die Fehlermeldung \texttt{testing.Testing\$} im Frontend angezeigt wurde. Die Ursache dafür, war das lediglich der geworfene Fehler weiter gereicht wurde und nicht tiefer in das Objekt eingegangen wurde. In der überarbeiteten Version wird die gesamte Fehlermeldung, wie in Tabelle~\ref{tab:check-api-fehlerbehandlung} visualisiert umgesetzt. Dabei wird ein einheitliches JSON, bestehend aus \texttt{Fehlercode}, \texttt{Status}, \texttt{errorType} also in welche Kategorie der Fehler fällt und die \texttt{message} also die Nachricht, welche erklärt, warum es möglicherweise zu diesem Fehler gekommen ist. Dieses System sorgt für eine wesentlich klarere Aufschlüsselung des Fehlers und macht es dem Nutzer im Frontend wesentlich einfacher zu erkennen, so das Problem liegen könnte. Der Wert \texttt{errorType} beschreibt die grobe Fehlerkategorie. Die konkrete Unterscheidung einzelner Ausführungsfehler erfolgt zusätzlich über die interne Fallbehandlung und die zurückgegebene Nachricht.

\begin{table}[h]
	\caption{Fehlerbehandlung und Rückgabemodell der Check-API}
	\label{tab:check-api-fehlerbehandlung}
	\begin{tabular}{p{3.1cm} p{2.5cm} p{8.5cm}}
		\textbf{Fehlerfall} & \textbf{\texttt{errorType}} & \textbf{Rückgabe} \\
		
		Erfolgreiche Prüfung &
		-- &
		\texttt{200 OK} mit \verb|{"status": "success"}| \\
		
		Fehlender Request-Body &
		-- &
		\texttt{400 Bad Request} mit \verb|{"status": "error",|
			\verb| "message": "Missing code to check."}| \\
		
		Sicherheitsverstoß &
		\texttt{security} &
		\texttt{400 Bad Request} mit \verb|{"status": "error",|
			\verb| "errorType": "security",|
			\verb| "message": "Security Violation:|
			\verb| Forbidden keyword detected."}|\\
		
		Kompilierfehler &
		\texttt{compilation} &
		\texttt{400 Bad Request} mit \verb|{"status": "error",|
			\verb| "errorType": "compilation",|
			\verb| "message": "Compilation Error:..."}| \\
		
		Fehlgeschlagener öffentlicher Test &
		\texttt{execution} &
		\texttt{400 Bad Request} mit \verb|{"status": "error",|
			\verb| "errorType": "execution",|
			\verb| "message": "..."}| \\
		
		Fehlgeschlagener versteckter Test &
		\texttt{execution} &
		\texttt{400 Bad Request} mit \verb|{"status": "error",|
			\verb| "errorType": "execution",|
			\verb| "message": "Hidden Requirement not met:}| 
			\verb| Your solution is not generic enough|
			\verb| (e.g., hardcoded values detected)."|\\
		
		Timeout / Endlosschleife &
		\texttt{execution} &
		\texttt{400 Bad Request} mit \verb|{"status": "error",|
			\verb| "errorType": "execution",|
			\verb| "message": "Execution Error:| 
			\verb| Timeout (Infinite loop detected?)"}|\\
		
		Sonstiger Laufzeitfehler &
		\texttt{execution} &
		\texttt{400 Bad Request} mit \verb|{"status": "error",|
			\verb| "errorType": "execution",|
			\verb| "message": "..."}| \\
	\end{tabular}
\end{table}




